% Template:     Informe LaTeX
% Documento:    Manual del proyecto - Game of Life en CPU, CUDA y OpenCL
% Codificación: UTF-8
%
% This is the developer MANUAL for the project. It is a standalone document,
% separate from the graded assignment report (main.tex). Build it the same way:
%   cd report && latexmk -pdf manual.tex   (or two pdflatex passes)

% CREACIÓN DEL DOCUMENTO
\documentclass[
	english,
	letterpaper, oneside
]{article}

% INFORMACIÓN DEL DOCUMENTO
\def\documenttitle {Project Manual: Game of Life on CPU and GPU}
\def\documentsubtitle {Architecture, mechanism, and a line-by-line code walkthrough}
\def\documentsubject {Computación en GPU}

\def\documentauthor {Matías Saavedra}
\def\coursename {Computación en GPU}
\def\coursecode {CC7515-1}

\def\universityname {Universidad de Chile}
\def\universityfaculty {Facultad de Ciencias Físicas y Matemáticas}
\def\universitydepartment {Departamento de Ciencias de la Computación}
\def\universitydepartmentimage {departamentos/dcc}
\def\universitydepartmentimagecfg {height=1.57cm}
\def\universitylocation {Santiago de Chile}

% INTEGRANTES, PROFESORES Y FECHAS
\def\authortable {
	\begin{tabular}{ll}
		Author:
		& \begin{tabular}[t]{l}
			Matías Saavedra
		\end{tabular} \\
		& \\
		\multicolumn{2}{l}{Document: developer manual} \\
		\multicolumn{2}{l}{Last updated: \today} \\
		\multicolumn{2}{l}{\universitylocation}
	\end{tabular}
}

% IMPORTACIÓN DEL TEMPLATE
\input{template}

% INICIO DE PÁGINAS
\begin{document}

% PORTADA
\templatePortrait

% CONFIGURACIÓN DE PÁGINA Y ENCABEZADOS
\templatePagecfg

% RESUMEN O ABSTRACT
\begin{abstractd}
	This document is the developer manual for the project: three
	implementations of Conway's Game of Life --- a sequential CPU baseline,
	a CUDA backend and an OpenCL backend --- built to benchmark
	cells-evaluated-per-second and report the GPU speed-up over the CPU. It
	explains the architecture from first principles, states the single
	invariant that shapes every design decision (\emph{the three backends
	must compute byte-for-byte identical boards so that any difference in
	speed is the real result}), and then walks the source code function by
	function, with the C\texttt{++} files explained line by line. At the time
	of writing the CPU backend is complete and acts as the correctness
	oracle; the CUDA and OpenCL backends are scaffolded but not yet
	implemented. The headline metric is fixed throughout:
	$\text{Mcells/s} = (\text{rows}\times\text{cols}\times\text{gens}) /
	(t \times 10^{6})$, measured against a no-op renderer so that rendering
	and host\,$\leftrightarrow$\,device transfers never pollute the number.
\end{abstractd}

% TABLA DE CONTENIDOS - ÍNDICE
\templateIndex

% CONFIGURACIONES FINALES
\templateFinalcfg

% ======================= INICIO DEL DOCUMENTO =======================

\section{What this manual is, and how to read it}
\label{sec:intro}

This is the engineering manual for the codebase, not the graded report. The
graded report lives in \texttt{main.tex} and presents experiments and results;
this manual explains \emph{how the code works and why it is built the way it
is}. Read it top to bottom the first time: the early sections establish the one
constraint that everything else follows from, and the later sections assume you
already understand it.

The manual blends two styles. It is structured like an engineering report
(setup, mechanism, tables, explicit trade-offs) but written like a deep code
explanation (every non-obvious decision is traced to \emph{why it exists} and
\emph{what would break if it were different}). For the C\texttt{++} translation
units it descends to a line-by-line account; for the headers, which are mostly
declarations, it summarises and shows the parts that carry real logic.

\paragraph{Build.} The manual compiles with the same vendored template as the
report:

\begin{sourcecode}{bash}{Building this manual.}
cd report && latexmk -pdf manual.tex
# or: pdflatex manual.tex && pdflatex manual.tex   (second pass for the TOC)
\end{sourcecode}

\section{The problem, and the invariant it forces}
\label{sec:problem}

The assignment looks like a simulation task but is really a \emph{measurement}
task. The deliverable is a number and an explanation of that number: how many
cells per second each backend evaluates, and \emph{why} the GPU wins (or loses)
at a given grid size. Conway's Game of Life --- here the HighLife variant
\texttt{B36/S23} --- is only the workload. It is deliberately a good workload
for the question being asked: a nine-point stencil over one-byte cells is
arithmetically trivial and memory-bound, so the experiment isolates memory
bandwidth and launch/transfer overhead, which is exactly the interesting story
for a GPU report.

Every architectural decision in the codebase serves one goal:

\begin{quote}
	\emph{Run the identical computation on three backends and trust that any
	difference in output is a bug, so that any difference in speed is the real
	result.}
\end{quote}

Two hard requirements fall out of that sentence, and they explain the entire
design:

\begin{enumerate}
	\item \textbf{Bit-for-bit equivalence across backends.} If CPU and GPU
	      produce different boards, the cells/sec comparison is meaningless.
	      This forces a \emph{shared} cell layout, a \emph{shared} rule, and
	      \emph{shared} edge handling.
	\item \textbf{Measurement that is not polluted by I/O.} The headline metric
	      is compute throughput, so rendering and host\,$\leftrightarrow$\,device
	      transfers must be removable from the timed path.
\end{enumerate}

The rule variant is the one detail most likely to be implemented inconsistently,
so it is worth stating precisely. It is \emph{not} vanilla Conway:

\begin{itemize}
	\item a dead cell is born on \textbf{exactly 3 OR exactly 6} live neighbours;
	\item a live cell survives on 2 or 3 live neighbours;
	\item everything else dies or stays dead.
\end{itemize}

The ``birth on 6'' clause is the easy mistake. It must hold identically in CPU,
CUDA and OpenCL, which is why it is written exactly once (Section
\ref{sec:liferules}) and pinned by a dedicated test (Section \ref{sec:tests}).

\section{Architecture at a glance}
\label{sec:arch}

The codebase is a \textbf{strategy pattern}: a backend-agnostic core plus two
swappable interfaces. The only thing that genuinely differs between backends ---
the per-generation compute --- sits behind \texttt{ISimEngine}; output sits
behind \texttt{IRenderer}; the application owns the loop and wires one of each at
runtime from the parsed configuration.

\begin{sourcecode}{plaintext}{The strategy layout. Parenthesised backends are scaffolded, not yet implemented.}
Application (main.cpp)  -- owns the loop + Config/CLI
  |-- ISimEngine   <- CpuEngine | (CudaEngine) | (OpenCLEngine)
  |-- IRenderer    <- NullRenderer | TextRenderer | (GuiRenderer later)
  '-- Grid + Pattern  (shared, backend-agnostic state & seed data)
\end{sourcecode}

The dependency arrows point \emph{from} engines \emph{to} the core, never back ---
that one-directional dependency is what keeps the core backend-agnostic and lets a
single host-side harness drive all three engines. The link graph is:

\begin{sourcecode}{plaintext}{Library dependency graph (CMake targets).}
gol_core  (Grid, Config+CLI, patterns, RLE, interfaces)
   |-- gol_cpu     (CpuEngine; links Threads::Threads)
   |-- gol_render  (TextRenderer)
   '-- (gol_cuda, gol_opencl -- opt-in, not yet present)
gol  executable  =  gol_core + gol_cpu + gol_render
\end{sourcecode}

The UML class diagram in Figure \ref{fig:uml} renders these relationships
formally. The \texttt{Application} depends only on the two strategy interfaces;
every backend and renderer is a realization behind them, and the shared
\texttt{Grid} flows through all of them unchanged.

\insertimage{uml_architecture}{width=\linewidth}{UML class diagram of the
project (generated from \texttt{report/diagrams/architecture.mmd} with
mermaid-cli). The \texttt{Application} owns one \texttt{ISimEngine} and one
\texttt{IRenderer} --- the two strategy interfaces. \texttt{CpuEngine} and
\texttt{CudaEngine} realize \texttt{ISimEngine}; \texttt{OpenCLEngine} is marked
\emph{planned}. \texttt{NullRenderer}, \texttt{TextRenderer} and
\texttt{AnsiRenderer} realize \texttt{IRenderer}. Dashed arrows are dependencies
(e.g. \texttt{CpuEngine} and \texttt{CudaEngine} both use
\texttt{LifeRules::nextState}; \texttt{RleLoader} produces a \texttt{Pattern}
that stamps a \texttt{Grid}). \label{fig:uml}}

Table \ref{tab:files} is the file map. Header-only units carry no \texttt{.cpp}
and impose no link dependency.

\begin{table}[H]
	\centering
	\caption{File map and status. ``HO'' = header-only (logic lives entirely in
	the header). The ``\S'' column points to the section that explains it.}
	\label{tab:files}
	\begin{tabular}{@{}lll@{}}
		\toprule
		\textbf{File} & \textbf{Responsibility} & \textbf{\S} \\
		\midrule
		\texttt{include/gol/Grid.hpp}            & Board state; the shared layout (HO) & \ref{sec:grid} \\
		\texttt{include/gol/LifeRules.hpp}        & The variant rule, written once (HO) & \ref{sec:liferules} \\
		\texttt{include/gol/Timer.hpp}           & Monotonic host stopwatch (HO) & \ref{sec:timer} \\
		\texttt{include/gol/ISimEngine.hpp}      & Engine strategy interface (HO) & \ref{sec:interfaces} \\
		\texttt{include/gol/IRenderer.hpp}       & Renderer strategy interface (HO) & \ref{sec:interfaces} \\
		\texttt{include/gol/Config.hpp}          & Run configuration struct (HO) & \ref{sec:config} \\
		\texttt{include/gol/engines/CpuEngine.hpp} & CPU engine declaration & \ref{sec:cpu} \\
		\texttt{include/gol/engines/CudaEngine.hpp} & CUDA engine declaration & \ref{sec:cuda} \\
		\texttt{include/gol/engines/cuda/kernel.cuh} & CUDA launcher declaration & \ref{sec:cuda} \\
		\texttt{include/gol/render/NullRenderer.hpp} & No-op renderer (HO) & \ref{sec:render} \\
		\texttt{include/gol/render/TextRenderer.hpp} & Text renderer declaration & \ref{sec:render} \\
		\texttt{include/gol/render/AnsiRenderer.hpp} & In-place ANSI renderer decl. & \ref{sec:render} \\
		\texttt{include/gol/patterns/Pattern.hpp} & Sparse seed pattern (HO) & \ref{sec:patterns} \\
		\texttt{include/gol/patterns/RleLoader.hpp} & RLE loader declaration & \ref{sec:patterns} \\
		\texttt{src/core/Config.cpp}             & CLI parsing & \ref{sec:configcpp} \\
		\texttt{src/core/main.cpp}               & The loop and the metric & \ref{sec:main} \\
		\texttt{src/engines/cpu/CpuEngine.cpp}   & CPU engine (seq + parallel) & \ref{sec:cpu} \\
		\texttt{src/engines/cuda/CudaEngine.cu}  & CUDA engine implementation & \ref{sec:cuda} \\
		\texttt{src/engines/cuda/kernel.cu}      & CUDA step kernels & \ref{sec:cuda} \\
		\texttt{src/patterns/RleLoader.cpp}      & RLE decoding & \ref{sec:patterns} \\
		\texttt{src/patterns/Pattern.cpp}        & Stamping a pattern onto a grid & \ref{sec:patterns} \\
		\texttt{src/render/TextRenderer.cpp}     & ASCII frame dump & \ref{sec:render} \\
		\texttt{src/render/AnsiRenderer.cpp}     & In-place ANSI animation & \ref{sec:render} \\
		\texttt{CMakeLists.txt}                  & Build, GPU opt-in gating & \ref{sec:cmake} \\
		\texttt{tests/*.cpp}                     & Executable spec of the invariants & \ref{sec:tests} \\
		\bottomrule
	\end{tabular}
\end{table}

\section{The shared core}
\label{sec:core}

The core is everything that must be identical across backends, plus the glue
that makes the backends interchangeable.

\subsection{\texttt{Grid} --- the spine of cross-backend equivalence}
\label{sec:grid}

\texttt{Grid} is the single most important object in the project. It is a
\textbf{flat, row-major \texttt{std::vector<unsigned char>}} where cell $(x,y)$
lives at index $y\cdot\text{cols}+x$, with $0=$ dead, $1=$ alive.

\begin{sourcecode}{cpp}{include/gol/Grid.hpp --- the load-bearing parts.}
class Grid {
public:
  Grid(std::size_t rows, std::size_t cols, unsigned char fillValue = 0)
      : rows_(rows), cols_(cols), cells_(rows * cols, fillValue) {}

  // Unchecked access by (x, y): the hot path must stay branch-free.
  unsigned char& at(std::size_t x, std::size_t y) { return cells_[y * cols_ + x]; }
  unsigned char  at(std::size_t x, std::size_t y) const { return cells_[y * cols_ + x]; }

  unsigned char* data() { return cells_.data(); }            // for host<->device transfers

  // Deterministic seeding: same seed => same board across all backends.
  void randomize(std::uint64_t seed, double aliveProbability = 0.3) {
    std::mt19937_64 rng(seed);
    std::bernoulli_distribution alive(aliveProbability);
    for (auto& c : cells_) c = alive(rng) ? 1u : 0u;
  }

  // Bit-for-bit comparison: the equivalence-test oracle check.
  bool operator==(const Grid& other) const {
    return rows_ == other.rows_ && cols_ == other.cols_ && cells_ == other.cells_;
  }
private:
  std::size_t rows_, cols_;
  std::vector<unsigned char> cells_; // row-major, size rows_ * cols_
};
\end{sourcecode}

\paragraph{Why this exact representation.} Each choice is forced by the
equivalence requirement:

\begin{itemize}
	\item \textbf{Flat and contiguous}, because a GPU device buffer is a flat
	      array. A flat \texttt{unsigned char} buffer maps to \texttt{cudaMemcpy}
	      / \texttt{clEnqueueWriteBuffer} with \emph{zero} translation --- the
	      same bytes on host and device. A \texttt{vector<vector<>>} would need a
	      transform on every transfer, and the transform itself could introduce
	      divergence.
	\item \textbf{One byte per cell, not one bit.} Bit-packing would be
	      $8\times$ smaller and use less bandwidth, but it would turn the
	      neighbour read into a bit-extraction and make the three backends harder
	      to keep identical. The project accepts the $8\times$ memory cost to keep
	      the stencil a plain array read. This is also \emph{why} shared memory is
	      expected to help little on the GPU (Section \ref{sec:timing}): a
	      one-byte stencil already fits comfortably in cache.
	\item \textbf{\texttt{at(x,y)} is unchecked.} The simulation hot path
	      evaluates this billions of times; a bounds branch per access would be a
	      measurable tax on a memory-bound kernel. Bounds are the caller's
	      responsibility.
\end{itemize}

\paragraph{The invariant.} \texttt{Grid} owns no device memory and contains no
device code. It is pure host state. Engines own their own ping-pong buffers and
touch the \texttt{Grid} only through \texttt{data()} during
\texttt{upload()}/\texttt{download()}. \textbf{What would break otherwise:} if the
\texttt{Grid} held device pointers, the host harness and the equivalence test
could no longer drive all three engines through one type.

Two methods exist purely for equivalence. \texttt{randomize} uses a seeded
\texttt{std::mt19937\_64}, so ``same seed $\Rightarrow$ same starting board''
holds across backends --- without a deterministic seed you could not start CPU
and GPU from an identical state. \texttt{operator==} compares dimensions
\emph{and} the full byte vector; this is literally the oracle check the
cross-backend test will call.

\subsection{\texttt{LifeRules} --- the rule written once, compiled three times}
\label{sec:liferules}

The rule variant is written exactly once and shared by source between the CPU
and CUDA backends.

\begin{sourcecode}{cpp}{include/gol/LifeRules.hpp --- the single source of truth for the rule.}
#ifdef __CUDACC__
#define GOL_FN __host__ __device__ inline
#else
#define GOL_FN inline
#endif

namespace gol {
// VARIANT (NOT vanilla Conway):
//   dead cell born on EXACTLY 3 OR EXACTLY 6 neighbours;
//   live cell survives on 2 or 3; everything else dies.
GOL_FN unsigned char nextState(unsigned char alive, int neighbors) {
  if (alive) {
    return (neighbors == 2 || neighbors == 3) ? 1u : 0u; // survive
  }
  return (neighbors == 3 || neighbors == 6) ? 1u : 0u;   // birth (3 or 6)
}
} // namespace gol
\end{sourcecode}

\paragraph{The mechanism.} When the host C\texttt{++} compiler reads this file,
\texttt{nextState} is an ordinary \texttt{inline} function used by
\texttt{CpuEngine}. When \texttt{nvcc} reads the \emph{same} file it defines
\texttt{\_\_CUDACC\_\_}, so the macro expands to \texttt{\_\_host\_\_
\_\_device\_\_} and the CUDA kernel calls the byte-for-byte identical logic. CPU
and CUDA therefore \emph{cannot} diverge on the rule --- they share source.

\paragraph{The acknowledged exception.} OpenCL C cannot \texttt{\#include} a
C\texttt{++} header, so the plan is to mirror this function in \texttt{kernel.cl}
(read at runtime, kept diff-able against this header). The project deliberately
does \emph{not} chase zero duplication between CUDA and OpenCL: two duplicated
copies that are easy to eyeball-diff beat one clever abstraction that is hard to
follow. That is a stated trade-off, not an oversight.

\subsection{\texttt{Timer} --- monotonic host timing}
\label{sec:timer}

\begin{sourcecode}{cpp}{include/gol/Timer.hpp --- the host stopwatch.}
class Timer {
public:
  Timer() : start_(Clock::now()) {}
  void reset() { start_ = Clock::now(); }
  double elapsedMillis() const {
    return std::chrono::duration<double, std::milli>(Clock::now() - start_).count();
  }
private:
  using Clock = std::chrono::steady_clock;   // monotonic, not system_clock
  Clock::time_point start_;
};
\end{sourcecode}

It uses \texttt{steady\_clock}, not \texttt{system\_clock}, because it is
monotonic: it cannot jump backwards if the wall clock is adjusted mid-run, which
would corrupt a timing. This is the CPU backend's kernel-timing source; the GPU
backends will time with \texttt{cudaEvent\_t} / CL events instead.

\subsection{\texttt{ISimEngine} and \texttt{IRenderer} --- the two strategy interfaces}
\label{sec:interfaces}

\begin{sourcecode}{cpp}{include/gol/ISimEngine.hpp --- the engine contract.}
class ISimEngine {
public:
  virtual ~ISimEngine() = default;
  virtual void upload(const Grid& initial) = 0;   // host -> "current" device buffer
  virtual void step() = 0;                         // read current, write next, swap
  virtual void download(Grid& out) = 0;            // device -> host
  virtual double lastKernelMillis() const = 0;     // compute-only time of last step()
  virtual std::string name() const = 0;            // "cpu", "cuda", ...
};
\end{sourcecode}

The crucial part is the documented \textbf{ping-pong contract}: an engine keeps
two equally-sized buffers; \texttt{step()} reads the ``current'' buffer, writes
the ``next'' buffer, then swaps them. This guarantees \texttt{step()} advances
exactly one generation and never aliases its input and output.

\paragraph{Why the two-buffer rule is non-negotiable.} A stencil reads each
cell's neighbours. If \texttt{step()} wrote new values into the same buffer it is
still reading from, cells processed later in the sweep would see
\emph{next-generation} neighbours instead of \emph{current-generation} ones, and
the result would silently depend on traversal order. The two-buffer invariant
prevents that, and it is stated in the interface so every backend obeys it. GPU
engines do the swap on the device; only \texttt{upload}/\texttt{download} cross
the bus.

\begin{sourcecode}{cpp}{include/gol/IRenderer.hpp --- the output contract.}
class IRenderer {
public:
  virtual ~IRenderer() = default;
  virtual void render(const Grid& grid, std::uint64_t generation) = 0;
  virtual bool shouldClose() const { return false; } // interactive renderers override
  virtual bool staysOpen()  const { return false; }  // hold final frame until quit?
};
\end{sourcecode}

\texttt{shouldClose()} defaults to \texttt{false}; only an interactive GUI
renderer would override it to break the loop early. The key implementation is
\texttt{NullRenderer} (Section \ref{sec:render}), whose entire purpose is to
\emph{not exist} in the timed path.

\subsection{\texttt{Config} --- the run configuration}
\label{sec:config}

\begin{sourcecode}{cpp}{include/gol/Config.hpp --- the knobs (abridged).}
enum class EngineKind { Cpu, Cuda, OpenCL };
enum class RendererKind { Null, Text, Ansi };

struct Config {
  std::size_t rows = 256, cols = 256;
  std::uint64_t generations = 100;
  EngineKind engine = EngineKind::Cpu;
  RendererKind renderer = RendererKind::Null;
  unsigned threads = 1;        // 1=sequential, N=parallel, 0=all hardware cores
  int blockSize = 256;         // GPU: threads per block (swept in the report)
  bool useShared = false;      // GPU: shared/local-memory tiling
  bool wrap = false;           // false=bounded edges, true=toroidal
  std::uint64_t seed = 1;
  std::optional<std::string> rlePath;
};
Config parse(int argc, char** argv);
\end{sourcecode}

The GPU-only knobs (\texttt{blockSize}, \texttt{useShared}) live here and are
simply ignored by \texttt{CpuEngine}. That is deliberate: one \texttt{Config}
type, and a single binary whose flag surface is the \emph{union} over all
backends, can script the entire benchmark matrix (block sizes
$\{32,64,128,256\}$ $\times$ shared/no-shared) without recompilation.

\section{Configuration parsing --- \texttt{Config.cpp}}
\label{sec:configcpp}

This translation unit turns \texttt{argv} into a fully-resolved \texttt{Config}.
Three file-private helpers carry the logic.

\begin{sourcecode}{cpp}{src/core/Config.cpp --- value consumption and integer parsing.}
[[noreturn]] void usageAndExit(int code) { /* prints help */ std::exit(code); }

std::string nextValue(int argc, char** argv, int& i, const char* flag) {
  if (i + 1 >= argc) { std::cerr << "error: " << flag << " requires a value\n";
                       usageAndExit(2); }
  return argv[++i];                       // consume the value AND advance caller's i
}

unsigned long long toULL(const std::string& s, const char* flag) {
  try { return std::stoull(s); }
  catch (const std::exception&) {
    std::cerr << "error: " << flag << " expects an integer, got '" << s << "'\n";
    usageAndExit(2);
  }
}
\end{sourcecode}

\texttt{usageAndExit} is marked \texttt{[[noreturn]]}: it never returns (it ends
in \texttt{std::exit}), which lets call sites fall through after calling it
without the compiler complaining about a missing return value.

\texttt{nextValue} is the mechanism that makes the parse loop correct. It takes
the loop index \texttt{i} \textbf{by reference}; line 36 checks a token follows
the flag, and line 40 \emph{pre-increments} \texttt{i} --- this both returns the
value token and advances the caller's index past it, so the main \texttt{for}
loop does not re-examine a flag's argument as if it were another flag. Remove the
\texttt{\&} on \texttt{i} and the parser would try to interpret every value as a
flag.

\texttt{toULL} centralises integer parsing: \texttt{std::stoull} throws
\texttt{std::invalid\_argument} on non-numeric input and
\texttt{std::out\_of\_range} on overflow; both are caught to print a friendly
message and exit, instead of an uncaught exception aborting the program.

\begin{sourcecode}{cpp}{src/core/Config.cpp --- the dispatch loop (representative branches).}
Config parse(int argc, char** argv) {
  Config cfg;                                   // start from defaults
  for (int i = 1; i < argc; ++i) {              // skip argv[0]
    const std::string arg = argv[i];
    if (arg == "-h" || arg == "--help") usageAndExit(0);
    else if (arg == "-r" || arg == "--rows")
      cfg.rows = (std::size_t)toULL(nextValue(argc, argv, i, "--rows"), "--rows");
    else if (arg == "--wrap")  cfg.wrap = true;            // boolean flag, no value
    else if (arg == "--rle")   cfg.rlePath = nextValue(argc, argv, i, "--rle");
    else if (arg == "--engine") {
      const std::string v = nextValue(argc, argv, i, "--engine");
      if (v == "cpu") cfg.engine = EngineKind::Cpu;
      else if (v == "cuda") cfg.engine = EngineKind::Cuda;
      else if (v == "opencl") cfg.engine = EngineKind::OpenCL;
      else { std::cerr << "error: unknown engine '" << v << "'\n"; usageAndExit(2); }
    }
    // ... --cols --gens --threads --seed --renderer --block --shared ...
    else { std::cerr << "error: unknown option '" << arg << "'\n"; usageAndExit(2); }
  }
  return cfg;
}
\end{sourcecode}

Each value-taking branch uses the idiom \texttt{toULL(nextValue(...))}, which
consumes and parses in one expression. \texttt{--wrap}, \texttt{--shared} are
value-less booleans. \texttt{--block}/\texttt{--shared} are accepted even in a
CPU-only build --- that is the union-flag-surface point in code.

\paragraph{What can go wrong.} Missing value, non-numeric value, unknown flag,
and unknown engine/renderer name are all caught. One unguarded edge: a negative
number passed to \texttt{--rows} is accepted by \texttt{stoull} as a huge
unsigned value (a known \texttt{stoull} quirk); it would surface immediately as a
failed allocation.

\section{The application loop and the measurement invariant --- \texttt{main.cpp}}
\label{sec:main}

\texttt{main.cpp} owns the loop, wires the chosen engine and renderer, seeds the
board, runs the generations, and prints the metric. It is where requirement~2
from Section \ref{sec:problem} --- keeping I/O out of the number --- is enforced.

\begin{sourcecode}{cpp}{src/core/main.cpp --- the factory and the seeder.}
std::unique_ptr<IRenderer> makeRenderer(RendererKind kind) {
  switch (kind) {
    case RendererKind::Text: return std::make_unique<TextRenderer>();
    case RendererKind::Ansi: return std::make_unique<AnsiRenderer>();
    case RendererKind::Null:
    default: return std::make_unique<NullRenderer>();   // safe default = benchmark-safe
  }
}

void seed(Grid& grid, const Config& cfg) {
  if (cfg.rlePath) {
    Pattern p = RleLoader::load(*cfg.rlePath);
    const std::size_t ox = cfg.cols > p.width  ? (cfg.cols - p.width)  / 2 : 0;
    const std::size_t oy = cfg.rows > p.height ? (cfg.rows - p.height) / 2 : 0;
    p.applyTo(grid, ox, oy);                              // centre the pattern
  } else {
    grid.randomize(cfg.seed);                             // deterministic fill
  }
}
\end{sourcecode}

The factory's \texttt{default} falling to \texttt{NullRenderer} is a safety
choice: any unexpected enum value degrades to the benchmark-safe no-op renderer
rather than crashing. In \texttt{seed}, the ternaries on the origins are
\emph{not} cosmetic: \texttt{cols}, \texttt{width} are unsigned, so if the
pattern were wider than the grid, \texttt{cfg.cols - p.width} would underflow to
a gigantic value and \texttt{applyTo} would clip the whole pattern off-grid. The
guard makes the origin $0$ in that case so at least the top-left lands.

\begin{sourcecode}{cpp}{src/core/main.cpp --- the loop and the metric.}
engine.upload(grid);

// Generation 0 is the seed itself: render it before the first step so frames
// are labelled seed, step 1, step 2, ... The NullRenderer path skips this.
if (rendering) renderer->render(grid, 0);

Timer wall;
double kernelMs = 0.0;
for (std::uint64_t gen = 0; gen < cfg.generations; ++gen) {
  engine.step();
  kernelMs += engine.lastKernelMillis();
  if (rendering) {                       // <-- the whole I/O path is gated here
    engine.download(grid);
    renderer->render(grid, gen + 1);     // state after gen+1 steps
    if (renderer->shouldClose()) break;
  }
}
const double wallMs = wall.elapsedMillis();

// Interactive viewers hold on the final frame until the user quits; staysOpen()
// is false for Null/Text, so this is skipped (and never hangs a pipe).
if (rendering && renderer->staysOpen())
  while (!renderer->shouldClose()) renderer->render(grid, cfg.generations);
renderer.reset(); // restore the terminal before printing the summary

const double cells = (double)cfg.rows * (double)cfg.cols * (double)cfg.generations;
auto mcellsPerSec = [cells](double ms) {
  return ms > 0.0 ? cells / (ms / 1000.0) / 1e6 : 0.0;
};
\end{sourcecode}

\paragraph{The measurement invariant, in code.} The \texttt{rendering} bool
(set to \texttt{cfg.renderer != RendererKind::Null}) gates the entire
\texttt{download} + \texttt{render} block. Under \texttt{NullRenderer} the loop
body is \emph{only} \texttt{step()} plus the time accumulation: no transfers, no
formatting. That is why the rule ``always benchmark against \texttt{NullRenderer}''
holds --- it is structurally enforced, not a convention you have to remember.

\paragraph{Two clocks.} \texttt{kernelMs} accumulates each step's
backend-native compute time (\texttt{lastKernelMillis}); \texttt{wall} measures
the whole loop. Reporting both separately lets the report distinguish ``the
kernel got faster'' from ``we are transfer-bound''.

\paragraph{Two subtleties.} The seed is generation $0$ and is rendered
\emph{before} the first \texttt{step()} so frame labels are honest (this is the
fix in commit \texttt{0f87fd9}); the \texttt{(double)} casts on the
\texttt{cells} product avoid 32-bit overflow at large grids (e.g.
$4096^2\times1000\approx 6.9\times10^{10}$ overflows 32-bit but is exact in
\texttt{double}), and the \texttt{ms > 0.0} guard avoids a divide-by-zero
yielding \texttt{inf}/NaN.

The whole body runs inside a \texttt{try/catch} so that, e.g., a missing RLE file
(which throws in the loader) is reported cleanly on \texttt{stderr} with exit
code 1. A guard before the loop rejects \texttt{--engine cuda|opencl} because
those backends are not built yet.

\section{The CPU engine --- \texttt{CpuEngine.cpp}}
\label{sec:cpu}

This is the most subtle file in the project, and the most important: it is the
correctness oracle for every future backend, and it implements the sequential
baseline and the data-parallel version \emph{in one code path}. The reason one
path can serve both is that the Game of Life step has \textbf{no
intra-generation data dependencies}: every cell reads only the current buffer and
writes only the next, so parallelising is \emph{purely} a matter of splitting the
rows across threads. The per-cell math is identical; only the row range differs.

\subsection{Data model and invariants}

\begin{sourcecode}{cpp}{include/gol/engines/CpuEngine.hpp --- the state (abridged).}
class CpuEngine final : public ISimEngine {
  // ... interface methods ...
  unsigned threads_;  bool wrap_;
  std::size_t rows_ = 0, cols_ = 0;
  std::vector<unsigned char> cur_;   // current generation (read by step)
  std::vector<unsigned char> nxt_;   // scratch for the next generation
  double lastMs_ = 0.0;
  // Persistent worker pool (only populated when threads_ >= 2):
  std::vector<std::thread> workers_;
  std::unique_ptr<std::barrier<>> barrier_;
  const unsigned char* src_ = nullptr;   // buffers published to workers each step
  unsigned char* dst_ = nullptr;
  std::atomic<bool> stop_{false};
};
\end{sourcecode}

\begin{itemize}
	\item \texttt{cur\_} / \texttt{nxt\_}: the two ping-pong buffers. Invariant:
	      at the start of every \texttt{step()}, \texttt{cur\_} holds the current
	      generation; at the end, after a swap, it holds the new one.
	\item \texttt{workers\_}: a \emph{persistent} thread pool, created only when
	      \texttt{threads\_ >= 2}. Persistent so per-generation thread
	      \emph{creation} never enters the timed region.
	\item \texttt{barrier\_}: a \texttt{std::barrier} with exactly
	      \texttt{threads\_} participants (main thread + \texttt{threads\_-1}
	      workers).
	\item \texttt{src\_} / \texttt{dst\_}: raw pointers ``published'' to the
	      workers each step so they know which buffers to read/write.
	\item \texttt{stop\_}: an atomic flag used only at destruction.
\end{itemize}

The class is non-copyable and non-movable: it owns \texttt{std::thread}s and a
\texttt{std::barrier}, neither of which has sensible copy/move semantics.

\subsection{Thread resolution, construction, and shutdown}

\begin{sourcecode}{cpp}{src/engines/cpu/CpuEngine.cpp --- lifecycle.}
unsigned CpuEngine::resolveThreads(unsigned requested) {
  if (requested != 0) return requested;
  unsigned hw = std::thread::hardware_concurrency();
  return hw == 0 ? 1u : hw;            // hardware_concurrency may report 0
}

CpuEngine::CpuEngine(unsigned threads, bool wrap)
    : threads_(resolveThreads(threads)), wrap_(wrap) {
  if (threads_ >= 2) startPool();      // sequential path touches no threads at all
}

CpuEngine::~CpuEngine() {
  if (!workers_.empty()) {
    stop_.store(true, std::memory_order_relaxed);
    barrier_->arrive_and_wait();       // release parked workers ONE last time
    for (auto& w : workers_) w.join();
  }
}
\end{sourcecode}

\texttt{resolveThreads} maps the request to a concrete count. \texttt{0} means
``all cores'', but \texttt{hardware\_concurrency()} is allowed to return $0$ when
it cannot detect the count, so it falls back to $1$. This guarantees
\texttt{threads\_ >= 1} always, which everything else relies on (e.g.
\texttt{rangeFor} divides by it).

The destructor is the trickiest correctness point in the file. At rest, every
worker is parked at the phase-1 barrier inside \texttt{workerLoop}. To shut down:
set \texttt{stop\_} (relaxed ordering suffices --- the barrier provides the
happens-before edge that publishes the flag), then \texttt{arrive\_and\_wait()}
\emph{once}. That single arrival completes phase 1 and releases every worker,
which then sees \texttt{stop\_} and \textbf{returns without arriving at the
phase-2 barrier}. That asymmetry is essential: if the workers hit phase 2 on
shutdown, the main thread (which is not waiting on phase 2 here) would leave them
blocked and \texttt{join()} would deadlock.

\subsection{The worker pool and the two-phase handshake}

\begin{sourcecode}{cpp}{src/engines/cpu/CpuEngine.cpp --- pool startup and worker loop.}
void CpuEngine::startPool() {
  // participants = main thread + (threads_ - 1) workers = threads_
  barrier_ = std::make_unique<std::barrier<>>((std::ptrdiff_t)threads_);
  workers_.reserve(threads_ - 1);
  for (unsigned id = 1; id < threads_; ++id)        // id 0 is the main thread
    workers_.emplace_back([this, id] { workerLoop(id); });
}

void CpuEngine::workerLoop(unsigned id) {
  while (true) {
    barrier_->arrive_and_wait();                     // phase 1: wait for work
    if (stop_.load(std::memory_order_relaxed)) return;
    auto [yBegin, yEnd] = rangeFor(id);
    computeRows(src_, dst_, yBegin, yEnd);
    barrier_->arrive_and_wait();                     // phase 2: signal completion
  }
}
\end{sourcecode}

Worker ids run $1\ldots\text{threads\_}-1$; \textbf{id 0 is the main thread}. The
main thread doing real work (partition 0) rather than just coordinating is a
deliberate efficiency choice: with $N$ cores you get $N$ workers, not $N-1$ plus
an idle coordinator.

The worker is a two-phase handshake. \textbf{Phase 1} (``go'') blocks until
\texttt{step()} or the destructor arrives. After release the worker checks
\texttt{stop\_} --- the only way the loop ever ends. Otherwise it computes its row
slice and signals \textbf{phase 2} (``done''), then loops back to park at phase 1.
The two distinct barriers per cycle separate ``all threads have started this
generation'' from ``all threads have finished it''.

\paragraph{Why the lock-free pointer publish is safe.} \texttt{step()} writes
\texttt{src\_}/\texttt{dst\_} \emph{before} arriving at phase 1; the worker reads
them \emph{after} leaving phase 1. The barrier establishes a happens-before
relationship, so the writes are visible without any lock or atomic on the
pointers themselves.

\subsection{Row partitioning --- \texttt{rangeFor}}

\begin{sourcecode}{cpp}{src/engines/cpu/CpuEngine.cpp --- the half-open row range for a partition.}
std::pair<std::size_t, std::size_t> CpuEngine::rangeFor(unsigned id) const {
  const std::size_t base = rows_ / threads_;
  const std::size_t rem  = rows_ % threads_;
  const std::size_t begin = id * base + std::min<std::size_t>(id, rem);
  const std::size_t count = base + (id < rem ? 1u : 0u);
  return {begin, begin + count};
}
\end{sourcecode}

The first \texttt{rem} partitions each get one extra row. Worked example with
\texttt{rows\_ = 17}, \texttt{threads\_ = 4} ($\text{base}=4$, $\text{rem}=1$):

\begin{table}[H]
	\centering
	\caption{\texttt{rangeFor} over 17 rows, 4 threads. Rows are contiguous,
	with no gaps or overlaps, and sum to 17.}
	\label{tab:rangefor}
	\begin{tabular}{@{}ccccc@{}}
		\toprule
		\texttt{id} & \texttt{base} & extra & \texttt{begin} & range \\
		\midrule
		0 & 4 & 1 & 0  & $[0,5)$ \\
		1 & 4 & 0 & 5  & $[5,9)$ \\
		2 & 4 & 0 & 9  & $[9,13)$ \\
		3 & 4 & 0 & 13 & $[13,17)$ \\
		\bottomrule
	\end{tabular}
\end{table}

The \texttt{min(id, rem)} term is the clever bit: partitions before \texttt{rem}
each contributed one extra row to the running offset, so partition \texttt{id}'s
start is shifted by \texttt{min(id, rem)}. Note that with \texttt{threads\_ == 1}
this returns $[0, \text{rows\_})$ --- the full grid --- so the sequential sweep is
the degenerate single-partition case. The \texttt{RemainderRowsHandled} test
(Section \ref{sec:tests}) targets exactly this logic.

\subsection{Neighbour counting and the stencil}

\begin{sourcecode}{cpp}{src/engines/cpu/CpuEngine.cpp --- both edge modes (condensed).}
int CpuEngine::countNeighbors(const unsigned char* g, std::size_t x,
                              std::size_t y) const {
  int n = 0;
  if (wrap_) {                                   // toroidal: indices wrap
    const long cols = (long)cols_, rows = (long)rows_;
    for (long dy = -1; dy <= 1; ++dy)
      for (long dx = -1; dx <= 1; ++dx) {
        if (dx == 0 && dy == 0) continue;        // a cell is not its own neighbour
        const std::size_t nx = ((long)x + dx + cols) % cols;   // +size before %
        const std::size_t ny = ((long)y + dy + rows) % rows;
        n += g[ny * cols_ + nx];
      }
  } else {                                       // bounded: OOB counts as dead
    for (long dy = -1; dy <= 1; ++dy)
      for (long dx = -1; dx <= 1; ++dx) {
        if (dx == 0 && dy == 0) continue;
        const long nx = (long)x + dx, ny = (long)y + dy;
        if (nx < 0 || ny < 0 || nx >= (long)cols_ || ny >= (long)rows_) continue;
        n += g[(std::size_t)ny * cols_ + (std::size_t)nx];
      }
  }
  return n;
}

void CpuEngine::computeRows(const unsigned char* src, unsigned char* dst,
                            std::size_t yBegin, std::size_t yEnd) const {
  for (std::size_t y = yBegin; y < yEnd; ++y) {
    const std::size_t row = y * cols_;           // hoisted out of inner loop
    for (std::size_t x = 0; x < cols_; ++x)
      dst[row + x] = nextState(src[row + x], countNeighbors(src, x, y));
  }
}
\end{sourcecode}

Cells are $0/1$, so \emph{summing} neighbour values \emph{is} counting the live
ones --- no comparison needed. In the toroidal branch the \texttt{+ size} before
\texttt{\%} is necessary because C\texttt{++}'s \texttt{\%} on a negative operand
yields a negative result; adding the size first guarantees a non-negative
dividend, so $-1 \bmod \text{cols}$ becomes $\text{cols}-1$ (wrap to the far
edge). In the bounded branch, out-of-range neighbours are simply skipped, i.e.
counted as dead --- this is the \textbf{default} edge mode. Both modes must match
every backend bit-for-bit, which is why both are tested.

\texttt{computeRows} is the single implementation of the stencil, called by
\emph{both} the sequential path (full range) and every worker (its slice). One
implementation of the math is the source of the sequential/parallel equivalence.
Hoisting \texttt{row = y*cols\_} out of the inner loop saves a multiply per cell.

\subsection{\texttt{upload}, \texttt{step}, \texttt{download}}

\begin{sourcecode}{cpp}{src/engines/cpu/CpuEngine.cpp --- the per-generation core.}
void CpuEngine::upload(const Grid& initial) {
  rows_ = initial.rows();  cols_ = initial.cols();
  cur_.assign(initial.data(), initial.data() + initial.size());
  nxt_.assign(cur_.size(), 0u);
}

void CpuEngine::step() {
  Timer t;                                            // CPU's lastKernelMillis source
  if (workers_.empty()) {
    computeRows(cur_.data(), nxt_.data(), 0, rows_);  // SEQUENTIAL: one full sweep
  } else {
    src_ = cur_.data();  dst_ = nxt_.data();          // publish buffers to workers
    barrier_->arrive_and_wait();                      // phase 1: release workers
    auto [yBegin, yEnd] = rangeFor(0);
    computeRows(src_, dst_, yBegin, yEnd);            // main thread = partition 0
    barrier_->arrive_and_wait();                      // phase 2: wait for all
  }
  std::swap(cur_, nxt_);                              // the ping-pong
  lastMs_ = t.elapsedMillis();
}

void CpuEngine::download(Grid& out) {
  std::copy(cur_.begin(), cur_.end(), out.data());   // host copy; device->host on GPU
}
\end{sourcecode}

\texttt{upload} records the dimensions, copies the seed into \texttt{cur\_}, and
sizes \texttt{nxt\_} (its initial contents do not matter --- \texttt{step}
overwrites every cell).

\texttt{step} is timed by the chrono \texttt{Timer}; this is the CPU's
\texttt{lastKernelMillis()} source. In the \textbf{sequential} branch it is one
\texttt{computeRows} over $[0,\text{rows\_})$ with zero synchronisation. In the
\textbf{parallel} branch the main thread mirrors the worker handshake from its
own side: publish the buffer pointers, phase-1 arrival releases the workers, the
main thread computes partition 0 itself instead of idling, and phase-2 arrival
blocks until every worker has finished. The symmetry is the point: \texttt{step}
runs the same phase1$\rightarrow$compute$\rightarrow$phase2 dance that each
worker runs, which is why they stay in lockstep generation after generation. The
\texttt{std::swap} is the ping-pong: \texttt{nxt\_} becomes the new
\texttt{cur\_}.

\texttt{download} copies the current buffer into the caller's grid. On the CPU
this is the ``no-op cost'' the interface mentions --- a host copy; on the GPU
backends this same method will be the device$\rightarrow$host transfer, which is
precisely why it is excluded from the benchmark.

\paragraph{What would break.} Remove the \texttt{stop\_} check and the destructor
deadlocks. Swap the order of ``publish \texttt{src\_/dst\_}'' and ``phase-1
arrive'' and workers could read stale pointers. Give \texttt{nextState} any
per-cell state and the sequential/parallel equivalence collapses --- it is
stateless on purpose.

\section{The CUDA engine --- \texttt{CudaEngine.cu} and \texttt{kernel.cu}}
\label{sec:cuda}

The GPU backend is the first that crosses the host/device boundary, yet it
reuses everything the core already fixed: the same flat \texttt{Grid} bytes, the
same \texttt{nextState} rule (now compiled \texttt{\_\_host\_\_ \_\_device\_\_}
under nvcc), and the same edge handling. That reuse is exactly why it can be
checked bit-for-bit against the CPU oracle (Section \ref{sec:tests}).

\subsection{Engine: device buffers, ping-pong, event timing}

\texttt{CudaEngine} owns two device buffers (\texttt{dCur\_}, \texttt{dNxt\_})
and does the ping-pong on the device; the host \texttt{Grid} is touched only by
\texttt{upload} (H2D \texttt{cudaMemcpy}) and \texttt{download} (D2H). The header
is deliberately CUDA-free (handles stored as \texttt{void*} /
\texttt{unsigned char*} and cast in the \texttt{.cu}), so \texttt{main.cpp}
compiles with the host compiler and only the \texttt{.cu} files go through nvcc.

\begin{sourcecode}{cpp}{src/engines/cuda/CudaEngine.cu --- one generation, kernel-only timing.}
void CudaEngine::step() {
  cudaEventRecord(start);
  launchLifeStep(dCur_, dNxt_, rows_, cols_, wrap_, blockSize_, useShared_);
  cudaEventRecord(stop);
  cudaEventSynchronize(stop);
  cudaEventElapsedTime(&ms, start, stop);  // kernel time only -- excludes H2D/D2H
  lastMs_ = ms;
  std::swap(dCur_, dNxt_);                  // ping-pong on the device
}
\end{sourcecode}

Bracketing the launch with \texttt{cudaEvent\_t} makes
\texttt{lastKernelMillis()} report kernel-only time, excluding the transfers ---
the GPU analogue of the CPU's chrono timing, and what makes the cells/sec numbers
comparable across backends (Section \ref{sec:metric}).

\subsection{Kernels: one thread per cell, 32-wide tiles}

\texttt{kernel.cu} holds two device kernels and a host launcher. The launcher
maps the block size onto a 2-D tile and selects the variant:

\begin{sourcecode}{cpp}{src/engines/cuda/kernel.cu --- launcher (tile mapping + kernel choice).}
const int bx = 32;                          // 32-wide rows -> coalesced global loads
const int by = (blockSize + bx - 1) / bx;   // 32/64/128/256 -> by = 1/2/4/8
dim3 block(bx, by);
dim3 grid((cols + bx - 1) / bx, (rows + by - 1) / by);
if (useShared)
  lifeShared<<<grid, block, (bx + 2) * (by + 2)>>>(src, dst, rows, cols, wrap);
else
  lifeGlobal<<<grid, block>>>(src, dst, rows, cols, wrap);
\end{sourcecode}

The 32-wide tile is the key performance choice: consecutive threads in a warp
read consecutive 1-byte cells of a row, so global loads coalesce --- the right
default for a memory-bound stencil. \texttt{lifeGlobal} reads the eight
neighbours straight from global memory. \texttt{lifeShared} (selected by
\texttt{--shared}) first stages a \texttt{(bx+2) x (by+2)} tile plus a one-cell
halo into shared memory cooperatively, applying the \emph{same} edge rule during
the halo load (bounded $\rightarrow$ 0, toroidal $\rightarrow$ wrap), then reads
neighbours from the tile. Both produce byte-identical results to the CPU oracle;
the report's job is to measure whether the staging pays off (Section
\ref{sec:timing} predicts it largely will not for a 1-byte stencil).

\subsection{The two swept options}

Block size (\texttt{--block}) and shared memory (\texttt{--shared}) are the two
assignment options chosen (\#1 and \#3); both are runtime-selectable, so one
binary drives the whole sweep. The CUDA target is opt-in
(\texttt{-DBUILD\_CUDA=ON}), built for \texttt{sm\_75} (the GTX 1660 Ti) with
\texttt{-lineinfo} so Nsight Compute can attribute time to source lines.

\section{The pattern pipeline --- \texttt{RleLoader.cpp} and \texttt{Pattern.cpp}}
\label{sec:patterns}

A \texttt{Pattern} is a \emph{sparse} seed: a bounding box plus the list of live
\texttt{(x,y)} offsets. Sparse because that is the natural output of RLE decoding
and it decouples a pattern from any particular grid size.

\begin{sourcecode}{cpp}{src/patterns/RleLoader.cpp --- the decode loop (core).}
std::size_t x = 0, y = 0, count = 0;
bool haveCount = false;
for (const char ch : body) {                       // body = payload, newlines dropped
  if (std::isspace((unsigned char)ch)) continue;
  if (std::isdigit((unsigned char)ch)) {
    count = count * 10 + (std::size_t)(ch - '0');  // build multi-digit run lengths
    haveCount = true; continue;
  }
  const std::size_t n = haveCount ? count : 1;     // count defaults to 1
  bool done = false;
  switch (ch) {
    case 'b': x += n; break;                                   // dead run: advance x
    case 'o': for (std::size_t i=0;i<n;++i) pat.liveCells.emplace_back(x++, y); break;
    case '$': y += n; x = 0; break;                           // end of row(s)
    case '!': done = true; break;                             // end of pattern
    default: break;                                            // ignore unexpected
  }
  count = 0; haveCount = false;
  if (done) break;
}
\end{sourcecode}

The RLE grammar is \texttt{<count><tag>} with the count defaulting to $1$. Digits
accumulate into \texttt{count}; on a tag, \texttt{n} is the accumulated count or
$1$, and the four tags drive a cursor: \texttt{b} skips dead cells, \texttt{o}
emits \texttt{n} live cells (advancing \texttt{x}), \texttt{\$} ends rows (and
\texttt{n\$} skips \texttt{n} blank rows at once --- how RLE compresses empty
space), \texttt{!} terminates. The payload is read into one \texttt{body} string
with newlines removed first, because RLE tokens can wrap across lines and cannot
be decoded line-by-line.

\paragraph{A concrete trace.} The \texttt{birth\_on\_six.rle} payload
\texttt{3o\$obo\$o!} decodes as:

\begin{sourcecode}{plaintext}{Decoding 3o\$obo\$o! -- six live cells around a DEAD centre (1,1).}
3o   -> (0,0) (1,0) (2,0)      x=3 y=0
$    -> y=1 x=0
o    -> (0,1)                  x=1
b    -> x=2                    (skips (1,1): the centre stays DEAD)
o    -> (2,1)                  x=3
$    -> y=2 x=0
o    -> (0,2)                  x=1
!    -> stop
\end{sourcecode}

That leaves a dead centre $(1,1)$ surrounded by exactly six live neighbours ---
which under the variant rule must be \emph{born}. This fixture is the payload
behind the ``birth on 6'' correctness test (Section \ref{sec:tests}).

Two robustness details round out the loader: the header parser
(\texttt{x = 3, y = 3, rule = B36/S23}) \emph{ignores} the \texttt{rule=} field
--- the simulation rule lives in code, not data, so a stray fixture cannot
silently change the simulation --- and if the header omits dimensions, the loader
infers the bounding box from $\max(x)+1$, $\max(y)+1$ over the decoded cells.
Opening a missing file throws \texttt{std::runtime\_error}.

\begin{sourcecode}{cpp}{src/patterns/Pattern.cpp --- stamping, with clipping.}
void Pattern::applyTo(Grid& grid, std::size_t originX, std::size_t originY) const {
  for (const auto& [dx, dy] : liveCells) {
    const std::size_t x = originX + dx, y = originY + dy;
    if (x < grid.cols() && y < grid.rows())   // clip, don't crash
      grid.at(x, y) = 1u;
  }
}
\end{sourcecode}

Because \texttt{x}, \texttt{y} are unsigned, the single test
\texttt{x < grid.cols()} rejects both ``too far right'' and any wrap-around at
once --- no separate \texttt{>= 0} check is needed. Cells that fall outside are
silently dropped, which is the right behaviour for a centering seeder: a pattern
slightly larger than the grid still produces a valid, cropped board rather than a
crash.

\section{Rendering --- \texttt{NullRenderer}, \texttt{TextRenderer}, \texttt{AnsiRenderer}}
\label{sec:render}

There are three renderers behind \texttt{IRenderer}, none of which is in the
benchmark path (the loop only renders when the renderer is not Null).
\texttt{NullRenderer::render} is an empty inline body --- header-only because
there is nothing to define. Its entire purpose is to \emph{not} exist in the
timed path (Section \ref{sec:main}).

\begin{sourcecode}{cpp}{src/render/TextRenderer.cpp --- one buffer, one write.}
void TextRenderer::render(const Grid& grid, std::uint64_t generation) {
  const std::size_t glyph = std::max(alive_.size(), dead_.size()); // bytes per cell
  std::string out;
  out.reserve(grid.size() * glyph + grid.rows() + 32);
  out += "generation "; out += std::to_string(generation); out += '\n';
  for (std::size_t y = 0; y < grid.rows(); ++y) {
    for (std::size_t x = 0; x < grid.cols(); ++x)
      out += (grid.at(x, y) ? alive_ : dead_);   // alive_/dead_ are std::string
    out += '\n';
  }
  out += '\n';
  std::fputs(out.c_str(), stdout);               // single write, no per-cell putchar
}
\end{sourcecode}

The mechanism worth noting is that it builds the entire frame into one
\texttt{std::string} (pre-\texttt{reserve}d so it never reallocates) and emits it
with a single \texttt{fputs}. Per-cell \texttt{putchar} would lock and flush the
stream thousands of times per frame. \texttt{at(x, y)} keeps the
\texttt{(column, row)} argument order matching \texttt{Grid}'s convention even
though storage is row-major. A terminal character cell is about twice as tall as
it is wide, so to make cells read as squares a live cell defaults to \emph{two}
full blocks (U+2588) and a dead cell to two spaces; the glyph fields are
\texttt{std::string} (not \texttt{char}) to hold the multi-byte, two-column cell,
and \texttt{alive\_} / \texttt{dead\_} share a display width so the grid stays
aligned. This renderer is for eyeballing correctness only --- never in a
benchmark.

\subsection{\texttt{AnsiRenderer} --- in-place animation, clipped to the viewport}
\label{sec:ansi}

\texttt{TextRenderer} \emph{appends} each frame, so the terminal scrolls and any
row wider than the terminal soft-wraps; it is a dump, not an animation.
\texttt{AnsiRenderer} instead homes the cursor and overwrites the same region
every frame using ANSI/DEC control sequences, producing a stable in-place
animation. It is a small-grid visualisation aid; benchmarks still use
\texttt{NullRenderer}.

\begin{sourcecode}{cpp}{src/render/AnsiRenderer.cpp --- the per-frame sequence (condensed).}
// Construction: hide cursor, disable autowrap (wide rows clip, not wrap), clear.
//   "\033[?25l"  "\033[?7l"  "\033[2J"
// Destruction (RAII): restore "\033[?7h" (wrap on) "\033[?25h" (cursor on).
void AnsiRenderer::render(const Grid& grid, std::uint64_t generation) {
  unsigned termRows = 24, termCols = 80;
  terminalSize(termRows, termCols);                      // ioctl(TIOCGWINSZ), 80x24 fallback
  const std::size_t visRows = std::min(grid.rows(), termRows > 1 ? termRows - 1u : 1u);
  const std::size_t cellCols = std::max<std::size_t>(1, displayWidth(alive_)); // cols/cell (2 by default)
  const std::size_t visCols = std::min(grid.cols(), termCols / cellCols);       // clip to viewport
  std::string out;
  out += "\033[?2026h";                                  // begin synchronized (atomic) update
  out += "\033[H";                                       // cursor home (overwrite, no flicker)
  out += "generation " + std::to_string(generation);
  if (visRows < grid.rows() || visCols < grid.cols())    // note the clip window
    out += " [" + std::to_string(visCols) + "x" + std::to_string(visRows) +
           " of " + std::to_string(grid.cols()) + "x" + std::to_string(grid.rows()) + "]";
  out += "\033[K\n";
  for (std::size_t y = 0; y < visRows; ++y) {
    for (std::size_t x = 0; x < visCols; ++x) out += (grid.at(x, y) ? alive_ : dead_);
    out += "\033[K\n";                                   // clear to EOL (wipe leftovers)
  }
  out += "\033[J";                                       // clear below (handles a shrunk terminal)
  out += "\033[?2026l";
  std::fputs(out.c_str(), stdout); std::fflush(stdout);
  if (delayMs_ > 0) std::this_thread::sleep_for(std::chrono::milliseconds(delayMs_));
}
\end{sourcecode}

\paragraph{Big grids: it clips, it does not wrap.} This is the key behavioural
difference. \texttt{TextRenderer} relies on the terminal: a $300\times300$ board
in an $80$-column terminal wraps every row and scrolls --- unreadable but
harmless. \texttt{AnsiRenderer} cannot allow wrapping, because a wrapped row
would occupy several terminal rows and desynchronise the cursor-home arithmetic
that the whole in-place scheme depends on. So it disables autowrap
(\texttt{\textbackslash 033[?7l}) and renders only the top-left
$\min(\text{rows},\,\text{term}-1)\times\min(\text{cols},\,\text{term}/\text{cell})$
window, where \emph{cell} is the per-cell display width ($2$ for the default
two-block glyph); it annotates the header, e.g. \texttt{generation 0 view(0,0)
40x23 of 300x300} ($300$ columns in an $80$-column terminal leaves room for $40$
two-column cells). The size is re-queried every frame, so resizing the terminal
is handled.

\paragraph{Interactive controls.} Because the window is smaller than the grid,
the renderer lets you drive it. When stdin is a TTY it puts the terminal in raw,
non-blocking mode (clearing \texttt{ICANON}/\texttt{ECHO}, \texttt{VMIN=VTIME=0})
and drains queued key presses each frame: the arrow keys shift the
\texttt{(offsetX\_, offsetY\_)} top-left of the viewport (by $\approx$ one eighth
of the view per press, clamped to the grid), \texttt{space}/\texttt{p} toggles
pause, and \texttt{q} sets the flag returned by \texttt{shouldClose()} so the main
loop stops. Reads are non-blocking, so input never stalls the simulation.

\emph{Pause is contained entirely in the renderer}: while paused,
\texttt{render()} loops on a \texttt{do/while} --- still polling input and
redrawing the \emph{same} generation (the header shows \texttt{[PAUSED]}) --- and
returns only once unpaused or quitting, so the application never calls
\texttt{step()} meanwhile. The main loop, the engine, and the \texttt{IRenderer}
interface are untouched, and the benchmark path (NullRenderer) never pauses. The
terminal is restored on destruction \emph{and} from a
\texttt{SIGINT}/\texttt{SIGTERM} handler (so Ctrl-C cannot leave it in raw mode);
when stdin is not a TTY, the controls are disabled and the top-left window is
shown.

\paragraph{Scrollbars.} To show \emph{where} the viewport sits, the renderer
draws a scrollbar in each clipped direction: a vertical bar in the rightmost
column (a heavy solid thumb, U+2503, over a dotted track, U+250A) and a
horizontal bar in a row under the grid (heavy U+2501 over dotted U+2508) --- the
solid thumb stands out against the dotted rail. Each is shown only when the grid
is clipped that way (one row/column is reserved for it). A shared helper
maps the visible band onto the bar --- the thumb length is the visible fraction
($\approx$ len\,$\times$\,visible\,/\,total) and its position tracks the offset
--- so the bar reads like any GUI scrollbar.

\paragraph{Holding after the run.} \texttt{staysOpen()} returns
\texttt{interactive\_}, so once the generation loop ends the application keeps
calling \texttt{render()} on the final board (still panning/pausing) until
\texttt{q}. Piped output reports \texttt{staysOpen() == false}, so a non-TTY run
exits normally and never hangs.

\paragraph{Robustness details.} A per-frame \texttt{\textbackslash 033[?2026h/l}
pair asks the terminal for a \emph{synchronized} (tear-free) update and is
silently ignored where unsupported; \texttt{\textbackslash 033[K} wipes any
leftover from a previously wider frame and \texttt{\textbackslash 033[J} wipes
rows orphaned by a shrunk terminal; the cursor is hidden during animation and
restored (with autowrap) by the destructor even on an early break or exception.
A small \texttt{delayMs} (default $50$\,ms) paces the otherwise-uncapped loop so
the animation is watchable; it lives only in this renderer, never in the timed
path. The size query is POSIX (\texttt{ioctl(TIOCGWINSZ)}); off a TTY it falls
back to $80\times24$.

\section{Build system --- \texttt{CMakeLists.txt}}
\label{sec:cmake}

The build must produce a working CPU\,+\,text binary on a machine with no CUDA
toolkit and no OpenCL. That requirement shapes the GPU targets.

\begin{sourcecode}{cmake}{CMakeLists.txt --- the core/CPU targets and the GPU double-gate.}
add_library(gol_core STATIC
    src/core/Config.cpp src/patterns/Pattern.cpp src/patterns/RleLoader.cpp)
target_include_directories(gol_core PUBLIC ${CMAKE_CURRENT_SOURCE_DIR}/include)

find_package(Threads REQUIRED)                       # std::thread / std::barrier
add_library(gol_cpu src/engines/cpu/CpuEngine.cpp)
target_link_libraries(gol_cpu PUBLIC gol_core Threads::Threads)

add_executable(gol src/core/main.cpp)
target_link_libraries(gol PRIVATE gol_core gol_cpu gol_render)

option(BUILD_CUDA "Build the CUDA engine" OFF)
if(BUILD_CUDA)
    set(_cuda_src ${CMAKE_CURRENT_SOURCE_DIR}/src/engines/cuda/CudaEngine.cu)
    if(EXISTS ${_cuda_src})                           # second gate: source must exist
        enable_language(CUDA)
        add_library(gol_cuda ${_cuda_src} .../kernel.cu)
        target_link_libraries(gol_cuda PUBLIC gol_core)
    else()
        message(WARNING "BUILD_CUDA=ON but source does not exist yet; skipping.")
    endif()
endif()
\end{sourcecode}

Each GPU backend is \textbf{double-gated}: an \texttt{option(... OFF)}
\emph{and} an \texttt{if(EXISTS <source>)} check. Even passing
\texttt{-DBUILD\_CUDA=ON} today emits a warning and skips the target because the
\texttt{.cu} source does not exist yet --- it does not fail configuration. The
OpenCL target follows the identical pattern. \texttt{find\_package(Threads
REQUIRED)} is what makes \texttt{std::thread}/\texttt{std::barrier} link on the
CPU engine.

\begin{sourcecode}{bash}{Building and testing.}
cmake -S . -B build                 # core + CPU + tests (GPU OFF by default)
cmake --build build
ctest --test-dir build --output-on-failure
# GPU opt-in (once the sources exist):
cmake -S . -B build -DBUILD_CUDA=ON -DBUILD_OPENCL=ON
\end{sourcecode}

The test \texttt{CMakeLists} prefers a system GoogleTest via
\texttt{find\_package(GTest CONFIG)} and otherwise \texttt{FetchContent}s
v1.17.0; it defines \texttt{PATTERNS\_DIR} so tests find the RLE fixtures
regardless of working directory, and registers each test with
\texttt{gtest\_discover\_tests}.

\section{Tests --- the executable specification of the invariants}
\label{sec:tests}

The tests are not an afterthought; they encode the invariants from Section
\ref{sec:problem} as runnable checks, in dependency order
(\texttt{CpuEngine} is the oracle, so it is verified first).

\begin{itemize}
	\item \textbf{\texttt{compile\_smoke\_test.cpp}} includes \emph{every}
	      interface header (so the scaffolding stays build-clean and mutually
	      consistent) but deliberately does not call the declared-but-undefined
	      symbols, so it links without those translation units. Its
	      \texttt{VariantRule} case pins the rule truth table directly.
	\item \textbf{\texttt{rules\_test.cpp}} checks \texttt{CpuEngine} on known
	      patterns: a block is a still-life, a blinker oscillates with period 2,
	      and --- the defining case --- a dead cell with exactly six live
	      neighbours is born.
	\item \textbf{\texttt{cpu\_parallel\_test.cpp}} proves the central design
	      claim. On a deliberately non-square, non-thread-divisible $128\times130$
	      grid it asserts thread counts $\{2,3,4,8,0\}$ all match the single-thread
	      reference \emph{bit-for-bit}, in both edge modes; a second case targets
	      \texttt{rangeFor}'s remainder logic with $17\times9$ and counts $4,5$.
	\item \textbf{\texttt{rle\_loader\_test.cpp}} covers parsing, dimensions,
	      specific coordinates (including that \texttt{birth\_on\_six}'s centre is
	      dead), the missing-file throw, and stamping with clipping.
\end{itemize}

\begin{sourcecode}{cpp}{tests/cpu\_parallel\_test.cpp --- the equivalence assertion.}
const Grid reference = runN(start, 1, wrap, 25);          // sequential oracle
for (unsigned threads : {2u, 3u, 4u, 8u, 0u /* all hw cores */}) {
  Grid parallel = runN(start, threads, wrap, 25);
  EXPECT_TRUE(parallel == reference)
      << "threads=" << threads << " wrap=" << wrap << " diverged";
}
\end{sourcecode}

\paragraph{The fourth test is still pending.} \emph{Cross-backend equivalence} ---
seed CPU/CUDA/OpenCL identically, run $N$ generations, assert bit-for-bit
equality against the \texttt{CpuEngine} oracle --- waits on the GPU engines. An
external oracle is also available: the headless \texttt{bgolly} runner simulates
\texttt{B36/S23} on an infinite grid, and the \texttt{highlife\_replicator.rle} /
\texttt{highlife\_spaceship.rle} fixtures already reproduce it bit-for-bit, so
they make good large-grid cross-checks for the GPU backends.

\section{Timing methodology and the expected GPU story}
\label{sec:timing}

\subsection{The metric}
\label{sec:metric}

The headline number is fixed across all three backends:

\begin{equation}
	\text{Mcells/s} =
	\frac{\text{rows}\times\text{cols}\times\text{generations}}
	     {t_{\text{elapsed}}\times 10^{6}}
\end{equation}

measured against \texttt{NullRenderer} so output and transfer cost never enter
$t$. Timing is uniform and in-program (not from a profiler) so the three backends
are measured the same way, exposed through
\texttt{ISimEngine::lastKernelMillis()}: \texttt{std::chrono} on CPU,
\texttt{cudaEvent\_t} around the launch on CUDA, and command-queue events with
\texttt{CL\_QUEUE\_PROFILING\_ENABLE} on OpenCL. Kernel time is reported
separately from host\,$\leftrightarrow$\,device transfers.

For the sweep data the binary has a CSV mode: \texttt{--csv} prints one data row
(columns \texttt{backend, rows, cols, generations, threads, block, shared, wrap,
kernel\_ms, wall\_ms, mcells\_kernel, mcells\_wall}) and \texttt{--csv-header}
prints just the header and exits, so a script writes the header once and appends
a row per configuration straight into Pandas. The speed-up is

\begin{equation}
	S(N) = \frac{T_{\text{CPU}}(N)}{T_{\text{GPU}}(N)}.
\end{equation}

\subsection{What the numbers should show, and why}

This is a memory-bound, arithmetically trivial nine-point stencil over one-byte
cells. That single fact predicts the shape of every curve the report will
produce:

\paragraph{Block size gives a concave curve that plateaus by 128--256.} Below
that, blocks are too small to hide memory latency (too few warps in flight to
cover the stalls of a bandwidth-bound kernel); past it, more threads per block do
not buy more effective bandwidth, so the curve flattens. The mechanism is
occupancy versus available memory throughput, not arithmetic.

\paragraph{Shared/local memory may give little or even negative gain.} The
classic stencil optimisation stages a tile plus a halo into shared memory to
reuse neighbour reads. Here each cell is one byte, so the working set of a tile
already lives in L1/L2; the explicit staging adds index arithmetic and a barrier
without saving meaningful DRAM traffic, and can lose. Explaining \emph{why} this
optimisation underperforms is one of the more interesting parts of the report.

\paragraph{No optimisation helps at small grids.} Kernel-launch and transfer
overhead dominate when there are few cells, so $S(N) < 1$ for small $N$ and the
CPU is competitive or faster. The sweeps must therefore run at large $N\times M$
to see the GPU win. Identifying the threshold $N$ where $S(N) > 1$ sustainably is
a graded deliverable.

The deep-dive profiling targets the \emph{best} parallel solution only:
\texttt{ncu} (Nsight Compute) for CUDA kernel-internal metrics --- occupancy,
stalls, memory throughput (it will not profile OpenCL) --- and \texttt{nsys}
(Nsight Systems) for the system timeline, which can also trace OpenCL on NVIDIA
hardware.

\section{Status, roadmap, and how to extend}
\label{sec:roadmap}

\paragraph{Done.} The backend-agnostic core, the CPU engine (sequential and
parallel in one path), both renderers, the RLE pipeline, the CLI, and the
\texttt{CudaEngine} (global + shared kernels, \texttt{--block}/\texttt{--shared},
cudaEvent timing), and the CSV benchmark-output mode
(\texttt{--csv}/\texttt{--csv-header}). Test suites: CPU rules,
sequential-vs-parallel, RLE loader, and the CUDA-vs-CPU cross-backend
equivalence (Section \ref{sec:cuda}). The CPU engine is the reference oracle.

\paragraph{Pending.} \texttt{OpenCLEngine} (\texttt{src/engines/opencl/}) and its
cross-backend test; the sweep scripts; and the benchmark runs, results CSVs, and
analysis notebooks that populate \texttt{main.tex} (including the Nsight deep-dive
on the best CUDA
config).

\paragraph{How a GPU backend slots in.} Implement \texttt{ISimEngine}: own two
device buffers, \texttt{upload} once, ping-pong on the device inside
\texttt{step}, \texttt{download} only when rendering. Reuse \texttt{nextState}
from \texttt{LifeRules.hpp} (CUDA includes it directly; OpenCL mirrors it in
\texttt{kernel.cl}). Match the edge handling of \texttt{countNeighbors} exactly,
in both bounded and toroidal modes. Wire it into \texttt{main.cpp}'s engine
selection and add the cross-backend test against the CPU oracle. Because the
\texttt{Grid} layout and the metric are already fixed, a correct new backend
needs no changes to the core, the harness, or the report's measurement code.

\section{Keeping this manual in sync}
\label{sec:maintenance}

This manual, the graded report (\texttt{main.tex}), and \texttt{CLAUDE.md} are
treated as living documents. The rule, recorded in \texttt{CLAUDE.md} under
``Documentation maintenance'', is: whenever a source file is added, removed, or
changed in a way that alters behaviour or structure, update all three in the
same change --- this manual's relevant section and file map (Table
\ref{tab:files}), the report if results or methodology move, and
\texttt{CLAUDE.md}'s ``Current state'' and layout. The most common trigger will
be landing the CUDA or OpenCL backend: that touches Sections \ref{sec:liferules},
\ref{sec:cpu} (as the oracle reference), \ref{sec:tests}, \ref{sec:timing}, and
\ref{sec:roadmap} here, plus the implementation and results sections of the
report.

% FIN DEL DOCUMENTO
\end{document}
